\label{ch:results}

This chapter presents, interprets, and summarises the real numerical
results of \texttt{notebooks/inference/inference4\_dayabay.ipynb}, executed
end to end against the real Daya Bay 8-detector, 8AD-period data
(\cref{ch:data-release}) using the physics of \cref{ch:detector-implementation}
and the inference machinery of \cref{ch:inference-process}. Every number
below is a real, reproducible output of that notebook, not illustrative.

\section{Real per-reactor relative weight}

\begin{table}[htb]
\centering
\caption{Real per-reactor relative power weight $w_r$ (\labelcref{eq:dayabay-relative-weight}), derived from the real weekly antineutrino-rate history, 8AD period.}
\label{tab:relative-weights-example}
\begin{tabular}{@{}lcccccc@{}}
\toprule
Reactor & R1 & R2 & R3 & R4 & R5 & R6 \\
\midrule
$w_r$ & 1.0010 & 0.9905 & 1.0126 & 0.9965 & 1.0079 & 0.9914 \\
\bottomrule
\end{tabular}
\end{table}

\Cref{tab:relative-weights-example} shows the real values obtained: a
small but genuine $\pm1$--$1.5\%$ real reactor-to-reactor spread around
the 6-reactor average of 1, replacing the previous implicit assumption
that all 6 real cores contribute identically to the flux sum
(\cref{eq:dayabay-face-flux}).

\section{Closure at Daya Bay's own published truth}
\label{sec:results-closure}

\begin{resultbox}
Evaluating the full real 8-detector model at Daya Bay's own published
$\theta_{13}$/$\Delta m^2_{31}$, with $\nu_{\rm norm}=1$ (no fit), gives
\begin{equation}
  N_{\rm pred}^{\rm total} = 3{,}705{,}407, \qquad
  N_{\rm obs}^{\rm total} = 3{,}114{,}934, \qquad
  \frac{N_{\rm pred}}{N_{\rm obs}} = 1.190,
\end{equation}
i.e.\ the simulation over-predicts the real observed rate by
$\approx19\%$, remarkably uniformly across all 8 detectors
($1.179$--$1.195$, \cref{tab:per-detector-ratio-truth}). Gradients
w.r.t.\ $\theta_{13}$, $\Delta m^2_{31}$, and $\nu_{\rm norm}$ are all
finite and non-zero, confirming autograd flows correctly through the real
multi-reactor flux, the real order-1/M cross section, and the real
IAV+LSNL+Gaussian response chain before any fit is attempted.
\end{resultbox}

\begin{table}[htb]
\centering
\caption{Predicted/observed ratio per detector at Daya Bay's own published truth, $\nu_{\rm norm}=1$.}
\label{tab:per-detector-ratio-truth}
\begin{tabular}{@{}lrrr@{}}
\toprule
Detector & Predicted & Real observed & Ratio \\
\midrule
AD11 & 827{,}463  & 695{,}290 & 1.190 \\
AD12 & 840{,}090  & 704{,}771 & 1.192 \\
AD21 & 792{,}729  & 667{,}823 & 1.187 \\
AD22 & 782{,}829  & 658{,}020 & 1.190 \\
AD31 & 116{,}030  &  97{,}069 & 1.195 \\
AD32 & 116{,}150  &  97{,}504 & 1.191 \\
AD33 & 115{,}008  &  96{,}793 & 1.188 \\
AD34 & 115{,}108  &  97{,}664 & 1.179 \\
\bottomrule
\end{tabular}
\end{table}

\section{Absolute-normalization audit}
\label{sec:results-audit}

A dedicated audit of the real $\sim19\%$ gap above checked every component
of the absolute-rate chain: the real target proton count and detector
efficiency (both loaded directly from the official release, not derived);
the order-1/M cross section (independently validated to $\sim1\%$ against
Vogel \& Beacom's own quoted coefficient, \cref{ch:detector-implementation});
the real Huber-Mueller antineutrino yield (the real table spans only
$1.75$--$13.0$~MeV, integrating to $\approx2.0$ antineutrinos/fission above
the IBD threshold -- the correct expected value for that truncated real
range, not a bug); and the thermal-power-to-fission-rate conversion (a
standard textbook formula). No unit or normalization error was found.

\begin{resultbox}
\Cref{tab:table5-crosscheck} cross-checks this notebook's real observed and
predicted rates-per-day against an \emph{independently published} table
(Table~5 of \cite{An2017DayaBayFluxSpectrum}, a different, earlier
621-day dataset). The real \emph{observed} rate matches the published table
to $1$--$3\%$ for every detector, confirming the real data/exposure loading
is correct; the \emph{predicted} rate exceeds it by the same
$\sim19$--$22\%$ already seen above, confirming the gap lives entirely in
the forward-model normalization chain, not in a data-loading or exposure
error.
\end{resultbox}

\begin{table}[htb]
\centering
\caption{Real observed/predicted rate per day vs.\ an independently published rate table \parencite{An2017DayaBayFluxSpectrum}.}
\label{tab:table5-crosscheck}
\begin{tabular}{@{}lrrrrr@{}}
\toprule
Detector & tpeanuts [1/day] & Real obs.\ [1/day] & Published [1/day] & obs/pub. & pred/pub. \\
\midrule
AD11 & 796.0 & 668.9 & 657.18 & 1.018 & 1.211 \\
AD12 & 810.7 & 680.1 & 670.14 & 1.015 & 1.210 \\
AD21 & 723.6 & 609.6 & 594.78 & 1.025 & 1.217 \\
AD22 & 714.9 & 600.9 & 590.81 & 1.017 & 1.210 \\
AD31 &  90.3 &  75.6 &  73.90 & 1.022 & 1.222 \\
AD32 &  90.4 &  75.9 &  74.49 & 1.019 & 1.214 \\
AD33 &  89.5 &  75.4 &  73.58 & 1.024 & 1.217 \\
AD34 &  89.6 &  76.0 &  75.15 & 1.012 & 1.193 \\
\bottomrule
\end{tabular}
\end{table}

Part of the real gap is a known, published physical effect: the
Huber-Mueller flux model itself is measured to predict $\approx5$--$6\%$
more than the real observed rate -- the ``Reactor Antineutrino Anomaly''
\parencite{An2017DayaBayFluxSpectrum}, quoted there as a measured/predicted
flux ratio of $0.946\pm0.020$ for the Huber+Mueller model. The remainder is
consistent with the many further real correlated systematics (target-mass,
efficiency, and flux-normalization nuisances) Daya Bay's own official
combined fit also carries and this project does not reproduce individually
-- precisely why the real official $\nu_{\rm norm}$ nuisance
(\cref{sec:nuisances}) exists rather than being fixed to 1.

\section{Full-spectral-shape fit: three attempts, one honest diagnosis}
\label{sec:results-fullshape}

\begin{figure}[htb]
\centering
\includegraphics[width=0.92\textwidth]{fig1_predicted_vs_real_spectra}
\caption{Predicted (tpeanuts, at Daya Bay's own published truth) vs.\ real
observed prompt-energy spectra, near-hall (AD11) and far-hall (AD31)
detectors.}
\label{fig:predicted-vs-real-spectra}
\end{figure}

\Cref{fig:predicted-vs-real-spectra} shows the real observed spectrum
against the \texttt{tpeanuts} simulation at Daya Bay's own published
truth for one near-hall and one far-hall detector: both follow the same
overall shape (rising then falling with a broad IBD-kinematics-driven peak
around 3--4~MeV), with the far-hall spectrum visibly suppressed relative
to its no-oscillation expectation, the real qualitative signature this
whole measurement is built on.

Minimising \labelcref{eq:poisson-nll} jointly over
$(\theta_{13},\Delta m^2_{31},\nu_{\rm norm})$ against the real
concatenated 8-detector, 208-bin spectrum was attempted three times, with
increasingly rich nuisance-parameter freedom, summarised in
\cref{tab:fullshape-attempts}.

\begin{table}[htb]
\centering
\caption{Full-spectral-shape fit: three real attempts (Section~5.1-5.3 of the notebook).}
\label{tab:fullshape-attempts}
\begin{tabular}{@{}lrrrr@{}}
\toprule
Attempt & $\hat\theta_{13}$ [deg] & $\Delta\hat m^2_{31}$ [$10^{-3}$eV$^2$] & $\hat\nu_{\rm norm}$ & $-2\ln L$ final \\
\midrule
Published truth & 8.506 & 2.528 & -- & -- \\
Unconstrained (5.1) & 9.172 & 7.483 & 0.874 & 6445.8 \\
+ Background nuisances (5.2) & 9.219 & 7.497 & 0.875 & 6390.6 \\
+ LSNL pull nuisances (5.3) & \textbf{0.000} & 11.872 & 0.819 & 4698.4 \\
\bottomrule
\end{tabular}
\end{table}

\begin{resultbox}
\textbf{None of the three attempts converges to a physical point.} The
unconstrained fit (5.1) already diverges: $\hat\theta_{13}=9.17^\circ$ ($+8\%$
absolute, $+15.7\%$ in $\sin^2 2\theta_{13}$) and
$\Delta\hat m^2_{31}=7.48\times10^{-3}$~eV$^2$, roughly $3\times$ the real
published value. Adding the 5 real Gaussian-constrained background-rate
nuisances (5.2, prior widths from \labelcref{eq:background-sigma}) improves
$-2\ln L$ only marginally (6390.6 vs.\ 6445.8, out of a $\sim90{,}000$-unit
total drop from the starting point) and leaves $\hat\theta_{13}$/$\Delta
\hat m^2_{31}$ essentially unchanged -- background-rate uncertainty is not
the dominant residual. Adding the 4 real LSNL pull-curve nuisances instead
(5.3) makes the fit \emph{worse}: $\hat\theta_{13}$ collapses to
$0^\circ$ rather than moving towards $8.51^\circ$, with two of the four
pulls at $3$--$5\sigma$ from their real unit-Gaussian prior. This is a
genuine degeneracy -- LSNL energy-scale freedom can disguise the
oscillation deficit itself as an energy-scale distortion when no
covariance structure constrains it jointly with the oscillation
parameters -- not evidence that LSNL uncertainty does not matter.
\end{resultbox}

\begin{notebox}
This three-way result is itself informative: it rules out both of the two
simplest candidate fixes (background normalization, LSNL curve freedom)
for the full-shape fit and points to the real, still-unimplemented next
step -- a proper correlated covariance or near/far-structured statistical
treatment (\cref{sec:results-summary}) -- rather than more isolated free
parameters.
\end{notebox}

\section{Rate-only fit}
\label{sec:results-rateonly}

\begin{figure}[htb]
\centering
\includegraphics[width=0.75\textwidth]{fig2_rateonly_fit_history}
\caption{Rate-only fit convergence: $-2\ln L$ (Poisson) at every LBFGS closure evaluation.}
\label{fig:rateonly-history}
\end{figure}

Integrating each detector's real spectrum to a single total count
(\labelcref{eq:rate-only-observable}) and fitting $\theta_{13}$/$\nu_{\rm
norm}$ jointly (with $\Delta m^2_{31}$ fixed at Daya Bay's own published
value) against the 8 real total counts:

\begin{resultbox}
\begin{equation}
  \hat\theta_{13} = 8.389^\circ \pm 0.139^\circ, \qquad
  \hat\nu_{\rm norm} = 0.8381\pm0.0008,
\end{equation}
$\sin^2 2\theta_{13}^{\rm fit}=0.08332$ vs.\ Daya Bay's real published
$0.08560$ -- a relative difference of $-2.66\%$, pull $-0.85\sigma$. Every
per-detector predicted/observed ratio at the fitted point lies within
$0.992$--$1.006$ of unity (\cref{fig:rateonly-2dscan}), a real, close
closure -- unlike \cref{sec:results-fullshape}.
\end{resultbox}

\begin{figure}[htb]
\centering
\includegraphics[width=0.62\textwidth]{fig3_rateonly_2d_scan}
\caption{2-D $-2\ln L$ scan over $(\theta_{13},\nu_{\rm norm})$, the rate-only model's full 2-parameter free set.}
\label{fig:rateonly-2dscan}
\end{figure}

\section{Near/far ratio fit: the headline result}
\label{sec:results-nearfar}

Fitting the real far-hall/near-hall per-bin count ratio
(\labelcref{eq:nearfar-ratio}) is the only strategy in this document with
genuine \emph{joint} sensitivity to both $\theta_{13}$ and $\Delta
m^2_{31}$ at once, and it is the only one that converges to a physical
point:

\begin{resultbox}
\begin{align}
  \hat\theta_{13} &= 8.445^\circ \pm 0.134^\circ,
  &\sin^2 2\theta_{13}^{\rm fit}&=0.08442, \\
  \Delta\hat m^2_{31} &= (2.4775\pm0.0548)\times10^{-3}\ {\rm eV}^2. &&
\end{align}
Relative to Daya Bay's own published values: $\sin^2 2\theta_{13}$
differs by $-1.38\%$, $\Delta m^2_{31}$ by $-2.01\%$ -- both within
$\sim1\sigma$. $-2\ln L$ falls from $114.3$ to $26.2$ over a real
26-bin ratio spectrum, a well-conditioned statistic (unlike the
full-shape fit's $\mathcal O(10^5)$ starting value, reflecting the
Poisson error on $\sim3\times10^6$ absolute counts).
\end{resultbox}

\begin{figure}[htb]
\centering
\includegraphics[width=0.85\textwidth]{fig6_nearfar_ratio_comparison}
\caption{Real observed far/near ratio vs.\ \texttt{tpeanuts} simulation (Daya Bay truth) and \texttt{tpeanuts} near/far fit.}
\label{fig:nearfar-ratio-comparison}
\end{figure}

\Cref{fig:nearfar-ratio-comparison} compares the real observed ratio
spectrum against the \texttt{tpeanuts} simulation at Daya Bay's own
published truth (no fit) and the \texttt{tpeanuts} near/far fit: even the
unfit truth-simulation curve already lies close to the real data (unlike
the absolute-count comparison of \cref{sec:results-closure}), because the
$\sim19\%$ normalization gap and much of the residual response mismatch
cancel by construction in the ratio.

\begin{figure}[htb]
\centering
\includegraphics[width=0.62\textwidth]{fig7_nearfar_2d_scan}
\caption{2-D $-2\ln L$ scan over $(\theta_{13},\Delta m^2_{31})$, the near/far model's full 2-parameter free set.}
\label{fig:nearfar-2dscan}
\end{figure}

\Cref{fig:nearfar-2dscan} shows the corresponding 2-D likelihood scan: the
\texttt{tpeanuts} fitted point and Daya Bay's own published point both sit
well inside the same low-$\Delta(-2\ln L)$ region, a direct visual
confirmation of the numerical agreement quoted above.

\begin{figure}[htb]
\centering
\includegraphics[width=0.75\textwidth]{fig5_nearfar_fit_history}
\caption{Near/far ratio fit convergence: $-2\ln L$ (chi2\_asymmetric) at every LBFGS closure evaluation.}
\label{fig:nearfar-history}
\end{figure}

\section{Real data comparison}

\begin{figure}[htb]
\centering
\includegraphics[width=0.85\textwidth]{fig4_real_comparison}
\caption{Real per-detector observed totals vs.\ \texttt{tpeanuts} simulation (Daya Bay truth) and \texttt{tpeanuts} rate-only fit.}
\label{fig:real-comparison}
\end{figure}

\Cref{fig:real-comparison} compares real per-detector observed totals
against the \texttt{tpeanuts} truth-simulation (visibly offset high by the
real $\sim19\%$ normalization gap, \cref{sec:results-audit}) and the
\texttt{tpeanuts} rate-only fit (visibly closer to the real data, having
absorbed that gap into $\hat\nu_{\rm norm}$).

\section{Summary of inference results and comparison with Daya Bay and NuFIT}
\label{sec:results-summary}

\Cref{tab:final-summary} collects every real fitted value obtained in this
chapter alongside Daya Bay's own published best fit
\parencite{An2023DayaBayPrecision} and, for $\theta_{13}$/$\Delta
m^2_{31}$, the independent global-fit value from NuFIT~6.1 (normal
ordering) \parencite{EstebanNuFit2024}, as used elsewhere in this project's
own preset registry (\code{tpeanuts.config.presets.\_SM\_NUFIT61\_NO}).

\begin{table}[htb]
\centering
\caption{Final comparison: tpeanuts inference vs.\ Daya Bay published vs.\ NuFIT 6.1.}
\label{tab:final-summary}
\begin{tabular}{@{}lrrr@{}}
\toprule
Quantity & tpeanuts (near/far, \S\ref{sec:results-nearfar}) & Daya Bay published & NuFIT 6.1$^\dagger$ \\
\midrule
$\sin^2 2\theta_{13}$ & $0.08442$ & $0.0856$ \parencite{An2023DayaBayPrecision} & $0.0878$ \\
$\theta_{13}$ [deg] & $8.445\pm0.134$ & $8.506$ & $8.62$ \\
$\Delta m^2_{31}$ [$10^{-3}$eV$^2$] & $2.478\pm0.055$ & $2.528$ & $2.511$ \\
\bottomrule
\end{tabular}
\vspace{2pt}

{\footnotesize $^\dagger$NuFIT 6.1's own global-fit best-fit central value
(normal ordering; itself informed by, among other inputs, Daya Bay's real
published measurement, not an independent cross-check) \parencite{EstebanNuFit2024}.}
\end{table}

\begin{resultbox}
\textbf{Final assessment.} The real near/far ratio fit
(\cref{sec:results-nearfar}) recovers both $\theta_{13}$ and $\Delta
m^2_{31}$ jointly, within $1$--$2\%$ of Daya Bay's own published values and
comfortably within NuFIT's global-fit range -- using \texttt{tpeanuts}'s
own real 6-reactor flux, real order-1/M cross section, real detector
response, and Daya Bay's own real near/far discovery strategy
\parencite{An2012DayaBayDiscovery}, not a synthetic or illustrative
dataset. The rate-only fit independently confirms $\theta_{13}$ alone to
similar precision. The full-spectral-shape fit, despite three real
attempts (including two different real nuisance-parameter families), does
not converge to a physical point, and this is reported as a genuine,
diagnosed finding rather than hidden: it demonstrates precisely why real
reactor-antineutrino collaborations build their combined analyses around a
full correlated systematic covariance, rather than a handful of
independent nuisance parameters, once the statistics are large enough that
sub-percent shape residuals dominate over the Poisson error.

\textbf{Recommended next step.} A real near/far \emph{spectral-shape}
covariance treatment (correlating the response/flux systematics between
the near and far halls directly, rather than only comparing their bin-wise
ratio) is the most promising avenue to make the full-spectral-shape fit
itself robust, now that both simpler nuisance-parameter fixes have been
tried and shown insufficient.
\end{resultbox}
